\section{Esercizi in preparazione all'esame}
\subsection{Esercizi SQL}
Si consideri la seguente base di dati:

\begin{itemize}
	\item \textbf{AUTO}(\underline{targa}, num\_telaio, modello, colore, km, cilindrata, posti, proprietario)
	\item \textbf{PRENOTAZIONE}(\underline{codice}, parcheggio, posto, auto, data\_prenotazione, data\_ora\_inizio, data\_ora\_fine, tipo, pagato)
	\item \textbf{PARCHEGGIO}(\underline{nome}, comune, indirizzo, coperto)
\end{itemize}

con i seguenti vincoli di integrità:

\begin{itemize}
	\item \texttt{PRENOTAZIONE.auto} $\rightarrow$ \texttt{AUTO}
	\item \texttt{PRENOTAZIONE.parcheggio} $\rightarrow$ \texttt{PARCHEGGIO}
\end{itemize}

\subsection{Esercizio 1}

Trovare tutti i parcheggi in cui vengono effettuate più prenotazioni di tipo online che prenotazioni fisiche.

\begin{verbatim}
	-- Numero di prenotazioni online
	SELECT P.nome
	FROM parcheggio P
	JOIN prenotazione R ON P.nome = R.parcheggio
	WHERE tipo = 'online'
	GROUP BY P.nome
	HAVING COUNT(*) > -- num prenotazioni online di un parcheggio
	(
	-- Numero di prenotazioni fisiche
	SELECT COUNT(*)
	FROM prenotazione R1
	WHERE R1.tipo = 'fisica'
	AND R1.nome = P.nome
	)
\end{verbatim}

La query esterna viene eseguita per ogni riga della query esterna. Usando il data
binding non si ha bisogno del \texttt{GROUP BY} perché si stanno selezionando le righe di uno
specifico parcheggio e quindi si vogliono contare tutte.

\subsection{Esercizio 2}

Trovare tutti i parcheggi in cui non sono mai state fatte prenotazioni fisiche nel 2024.

\subsubsection*{Soluzione 1 — EXCEPT}

\begin{verbatim}
	SELECT P.NOME
	FROM PARCHEGGIO P
	
	EXCEPT
	
	SELECT R.PARCHEGGIO
	FROM PRENOTAZIONE R
	WHERE R.DATA_PRENOTAZIONE >= '01/01/2024'
	AND R.DATA_PRENOTAZIONE <= '31/12/2024'
	AND R.TIPO = 'FISICA';
\end{verbatim}

\subsubsection*{Soluzione 2 --- NOT IN}

\begin{verbatim}
	SELECT P.NOME
	FROM PARCHEGGIO P
	WHERE P.NOME NOT IN (
	SELECT R.PARCHEGGIO
	FROM PRENOTAZIONE R
	WHERE EXTRACT(YEAR FROM R.DATA_PRENOTAZIONE) = 2024
	AND R.TIPO = 'FISICA'
	AND R.PARCHEGGIO IS NOT NULL
	);
\end{verbatim}
\subsection{Esercizio 3}

Trovare quanti parcheggi del comune di Verona hanno avuto un numero medio di prenotazioni complessive nel 2024 superiore al numero medio di prenotazioni dei parcheggi del comune di Padova.

\begin{verbatim}
	CREATE TEMP VIEW nprenotazioni AS (
	SELECT P.nome, P.comune, COUNT(*) AS num
	FROM parcheggio P
	JOIN prenotazioni R ON P.nome = R.parcheggio
	WHERE R.data_prenotazione >= '01/01/2024'
	AND R.data_prenotazione <= '31/12/2024'
	GROUP BY P.nome, P.comune
	);
	
	SELECT COUNT(*)
	FROM nprenotazioni N
	WHERE N.comune = 'Verona'
	AND N.num > (
	SELECT AVG(N1.num)
	FROM nprenotazioni N1
	WHERE N1.comune = 'Padova'
	);
\end{verbatim}